% ======================================================================
% Section 2 — Preliminaries
% ======================================================================

\section{Preliminaries}\label{sec:prelim}

\subsection{Notation}

Throughout, $n\ge 3$ is the number of speeds, $V=(v_1,\dots,v_n)$ has
distinct positive integer entries with $\gcd(V)=1$ and $v_1=1$ when a
coordinate model is in use, and $\wnc{x}=\min(x-\floor{x},
\floor{x}+1-x)$ is the distance to the nearest integer. We write
$\lambda(V)=\max_{t\in\T}\min_p \wnc{v_p t}$ for the optimal loneliness
and $M(V)$ for the same quantity; the two notations are interchangeable
and we keep $\lambda$ in metric contexts and $M$ in the
pair-sum context. The threshold is
$\mathrm{thr}_n=(n-1)/(2(n+1))=1/2-1/(n+1)$. Two speed families recur:
the \emph{harmonic} family $V=(1,2,\dots,n)$ and the \emph{rung-2}
family $V=(1,\dots,n-1,2n)$. For these, $\lambda$ is classical:
$\lambda(1,\dots,n)=1/(n+1)$ exactly (the harmonic family is extremal,
with zero margin against the conjecture), and
$\lambda(1,\dots,n-1,2n)=2/(2n+1)$ exactly (the rung-2 family is not
extremal; its margin is $1/((n+1)(2n+1))$). Both identities are
re-asserted exactly for $n\le 10$ in the provenance run of
Appendix~\ref{sec:prov}.

\subsection{Four forms of the covering statement}

Write $\delta(V)=\tfrac12-\lambda(V)$ and let $\mathbf{1}=(1,\dots,1)$.
The following forms of $\lrc_n$ are used interchangeably in the
literature and in this program, and it is essential to keep their
strengths apart.

\begin{enumerate}[leftmargin=1.6em,itemsep=2pt]
\item \textbf{Loneliness form.} $\lambda(V)\ge 1/(n+1)$ for all $V$.
\item \textbf{Center-depth form ($\delta$-form).}
$\delta(V)\le \mathrm{thr}_n$, where
$\delta(V)=\operatorname{dist}_\infty\!\bigl(\R v,\;\Z^n+\tfrac12\mathbf{1}\bigr)$
is the $\ell_\infty$ distance from the trajectory line to the shifted
lattice. The identity
$\delta=\tfrac12-\lambda$ is elementary: for every $t$,
$\max_j \wnc{\tfrac12-v_jt}=\tfrac12-\min_j\wnc{v_jt}$.
\item \textbf{Quotient-torus form.} Let $\mathcal{X}_v=\R^n/\R v$ be
the quotient vector space (the calligraphic symbol avoids colliding with
the speed vector $V$), let $\pi\colon\R^n\to\mathcal{X}_v$ be the
projection, $\Lambda=\pi(\Z^n)\cong\Z^{n-1}$, and let
$\|\cdot\|_{\mathcal{X}_v}$ be the quotient norm
$\|c\|_{\mathcal{X}_v}=\operatorname{dist}_\infty(x,\R v)$ for any lift
$x$ of $c$. Then $\delta(V)=\operatorname{dist}_{\mathcal{X}_v}(c^*,\Lambda)$
where $c^*=\pi(\tfrac12\mathbf{1})$ is the \emph{center class}, a
$2$-torsion point whose coordinates are $\tfrac12$ on even speeds and $0$
on odd speeds. The $\delta$-form reads: the center class has quotient-norm
depth at most $\mathrm{thr}_n$ in every $\mathcal{X}_v$.
\item \textbf{Covering-radius form ($\rho$-form).}
Let $\rho(V)=\max_{c\in\mathcal{X}_v}\operatorname{dist}_{\mathcal{X}_v}(c,
\Lambda)$ be the covering radius of $\Lambda$ in the quotient norm. The
$\rho$-form of the conjecture asserts $\rho(V)\le\mathrm{thr}_n$ for all
$V$.
\end{enumerate}

Forms (1)--(3) are equivalent to $\lrc_n$: this is the content of the
shift identity above together with the standard quotient-torus dictionary,
and the equivalence is asserted computationally, exactly, on every
instance of the program (the geometric pipeline is validated end to end
in the runs of Appendix~\ref{sec:prov}). Form (4) is \emph{strictly
stronger}. Since the maximum defining $\rho$ ranges over all classes and
$\delta$ is the depth of one class, $\delta(V)\le\rho(V)$ always; hence
$\rho\le\mathrm{thr}_n$ implies the $\delta$-form and thus $\lrc_n$. The
converse fails, and the failure is not an artifact of small examples: it
begins exactly at $n=5$ and persists (\S\ref{sec:tm-zono},
\S\ref{sec:witness}). At the harmonic instances the $\delta$-form is
tight ($\delta=\mathrm{thr}_n$), so the $\rho$-form has no slack to
absorb the difference between the center and the deepest hole: the
two statements coincide on all tested instances for $n\le 4$---more
precisely, the $\rho$-form \emph{holds} on every instance tested at
$n\le 4$, certified through $m=24$ at $n=3$ and $m=16$ at $n=4$, while
the center-is-deepest property holds exactly at the multiples---which is
why the identification was easy to believe---and separate from $n=5$ on,
with exact witnesses.

For computations we use two further standard facts, both proved and
machine-validated in the program's records \cite{prog-artifacts}. First,
in \emph{$c$-coordinates} $c_j=x_j-v_jx_1$ ($j=2,\dots,n$; valid when
$v_1=1$) the lattice is exactly $\Z^{n-1}$ and the quotient norm is the
maximum of the finitely many linear forms
\[
\phi^{1j}=\frac{c_j}{1+v_j},\qquad
\phi^{ij}=\frac{v_jc_i-v_ic_j}{v_i+v_j}\quad(2\le i<j),
\]
each of $\ell_1$-norm at most $1$; hence
$\|c\|_{\mathcal{X}_v}\le\|c\|_\infty$ and Lipschitz certification
needs no norm-equivalence constant. Second, the depth function admits the
one-parameter \emph{runner formulation}
\begin{equation}\label{eq:runnerform}
D(c)\;=\;\min_{\sigma\in[0,1)}\;\max\Bigl(\wnc{\sigma},\;
\max_{j\ge 2}\wnc{c_j-v_j\sigma}\Bigr),
\end{equation}
which is the distance of $c$ to $\Lambda$ on the quotient torus; this is
the form used by the branch-and-bound certifier of \S\ref{sec:witness},
and it was validated against the independent linear-form/lattice-search
implementation on random rational points for every instance used in this
paper (forty points per instance, exact agreement throughout).

\subsection{The Lonely Runner zonotope and its Ehrhart theory}

The \emph{LR-zonotope} of $V$ is $Z(V)=\pi([0,1]^n)\subset
\mathcal{X}_v$; it is an $(n{-}1)$-dimensional lattice zonotope with
respect to $\Lambda$, with relative volume $\sum_j v_j$ (normalized
volume $d!\,\sum_j v_j$, so that $h^*(1)=d!\,\sum_j v_j$)
\cite{Shephard,BeckRobins}. Its lattice-point geometry encodes the
covering problem above: holes of the $\Lambda$-covering correspond to
classes of large quotient depth. Its Ehrhart theory is arithmetic and
clean \cite{BeckRobins,dAMoci,StanleyEC}. Writing $d=n-1$, the Ehrhart
polynomial of $Z(V)$ is
\begin{equation}\label{eq:ehr}
\mathrm{Ehr}_{Z(V)}(t)\;=\;\sum_{k=0}^{d} e_k\,t^k,\qquad
e_k\;=\;\sum_{\substack{S\subseteq[n]\\|S|=k}}\gcd\!\bigl(v_{S^c}\bigr),
\qquad e_d=\sum_j v_j,
\end{equation}
and the $h^*$-polynomial is obtained
from the Eulerian expansion
$\sum_{t\ge0}\mathrm{Ehr}(t)z^t=h^*(z)/(1-z)^{d+1}$; explicitly
$h^*(z)=\sum_k e_k\,(1-z)^{d-k}A_k(z)$ with $A_k$ the Eulerian
polynomials \cite{Frobenius}. Formula \eqref{eq:ehr} was derived in the
program's records and verified in three independent ways on every
instance used here (the Eulerian expansion, the standard Ehrhart
inversion, and a depth-first coset count at $t=1,2$;
Appendix~\ref{sec:prov}). The gcd structure is that of the arithmetic
matroid of the projected basis \cite{dAMoci}, which is why the
arithmetic-matroid literature is the natural home for the residual
positive question in Appendix~\ref{sec:hstar}.

\subsection{The certification ladder}

Every classification claim in this paper carries one of three labels.
\emph{Exact}: a complete proof in rational arithmetic (the $d=2$
arrangement enumerations; all refuting witnesses, which are single
evaluations of $D$ at rational points). \emph{Certified}: a complete
proof in which an outer numerical loop (a grid or a branch-and-bound
frontier) is combined with an exact inner analysis and a rigorous
one-sided error bound (the $d=3$ away-certifications; the
branch-and-bound certificates of \S\ref{sec:witness}). \emph{Consistent}:
search-level evidence only (grid maxima equal to the claimed value,
multi-start optimization finding nothing above it), with no completeness
claim. The labels exist precisely because enclosure-level evidence must
not be presented as proof; the ladder is applied uniformly to every
claim in this paper.
